\section{\S 3.4 金属腐蚀与防护}\label{sect:金属腐蚀与防护}

    \subsection{金属腐蚀的本质}\label{subsec:金属腐蚀本质}

        \textcolor{red}{\cemh{M - $n$ e- \xleq{} M^{n+}}（金属被氧化——被腐蚀）}

    \subsection{化学腐蚀}\label{subsec:化学腐蚀}

        金属与氧化剂（\cemh{O2}、\cemh{Cl2}、\cemh{H+} 等）直接接触发生氧化还原反应。

        \textcolor{blue}{化学腐蚀较少且缓慢。}

    \subsection{电化学腐蚀}\label{subsec:电化学腐蚀}

        \textcolor{red}{电化学腐蚀是更普遍的金属腐蚀形式。}

        铁制品（铁碳合金）在电解质溶液中形成微电池。

        \subsubsection{吸氧腐蚀}\label{subsubsec:吸氧腐蚀}

            \textcolor{red}{更普遍的电化学腐蚀形式（中/碱性条件）。}

            负极（\cemh{Fe}）：\cemh{Fe - 2e- \xleq{} Fe^{2+}}

            正极（\cemh{C}）：\cemh{O2 + 2H2O + 4e- \xleq{} 4OH-}

            铁锈的生成：

            \[\cemh{Fe^{2+}} \xleq[]{\text{氧化}} \cemh{Fe(OH)2} \xleq[]{\text{氧化}} \cemh{Fe(OH)3} \xleq{} \cemh{Fe2O3*$n$H2O}\]（铁锈）

            \textcolor{blue}{\(\Delta S < 0\) 且自然条件下自发（\(\Delta H < 0\)）。}

        \subsubsection{析氢腐蚀}\label{subsubsec:析氢腐蚀}

            \textcolor{red}{酸性条件下的电化学腐蚀。}

            负极（\cemh{Fe}）：\cemh{Fe - 2e- \xleq{} Fe^{2+}}

            正极（\cemh{C}）：\cemh{2H+ + 2e- \xleq{} H2 ^}

    \subsection{金属腐蚀的利用}\label{subsec:金属腐蚀的利用}

        铁炭混合物（铁屑和活性炭的混合物）在水溶液中可形成许多微电池，用于去除水中 \cemh{Cu^{2+}}、\cemh{Pb^{2+}} 等重金属离子。

    \subsection{金属的防护}\label{subsec:金属的防护}

        \subsubsection{覆盖保护层}\label{subsubsec:覆盖保护层}

            \begin{itemize}
                \item 涂装（油漆）、电镀
                \item 马口铁——镀锡；白口铁——镀锌
            \end{itemize}

        \subsubsection{改变成分}\label{subsubsec:改变成分}

            减少含碳量；合金化（如不锈钢：\cemh{Fe} + \cemh{Cr} + \cemh{Ni}）。

        \subsubsection{电化学保护法}\label{subsec:电化学保护法}

            \begin{dfn}[牺牲阳极法]{牺牲阳极的阴极保护法}
                将活泼金属（\cemh{Zn}、\cemh{Mg}、\cemh{Al}）与被保护的钢铁设备连接，活泼金属作阳极被腐蚀，钢铁作阴极被保护。

                \textcolor{blue}{例如：轮船外壳绑锌块。}

                阳极（\cemh{Zn}）：\cemh{Zn - 2e- \xleq{} Zn^{2+}}

                阴极（\cemh{Fe}）：\cemh{O2 + 2H2O + 4e- \xleq{} 4OH-}
            \end{dfn}

            \begin{dfn}[外加电流法]{外加电流的阴极保护法}
                直流电源负极接被保护金属（阴极），辅助阳极（石墨/\cemh{Pt}/\cemh{Ti}）接电源正极。

                \textcolor{blue}{例如：埋地管道、大型钢结构。}

                \textcolor{red}{强迫阴极（原本的负极）得电子。}
            \end{dfn}

            \begin{center}
            \small
            \begin{tabular}{|c|c|c|}
                \hline
                \textbf{对比项} & \textbf{牺牲阳极法} & \textbf{外加电流法} \\
                \hline
                是否需要电源 & 不需要 & 需要 \\
                \hline
                适用范围 & 小结构 & 大型结构 \\
                \hline
                保护效果 & 寿命有限 & 可调节 \\
                \hline
            \end{tabular}
            \end{center}

        \subsubsection{环境调控}\label{subsubsec:环境调控}

            干燥环境、除去溶解氧、加入缓蚀剂等。
